\section{Minimal polynomials}
\label{section-minimal-polynomials}

\noindent
Let $E/k$ be a field extension, and let $\alpha \in E$ be algebraic over $k$.
Then $\alpha$ satisfies a (nontrivial) polynomial equation in $k[x]$.
Consider the set of polynomials $P \in k[x]$ such that $P(\alpha) = 0$; by
hypothesis, this set does not just contain the zero polynomial.
It is easy to see that this set is an {\it ideal.} Indeed, it is the kernel
of the map
$$
k[x] \to E, \quad x \mapsto \alpha
$$
Since $k[x]$ is a PID, there is a {\it generator} $P \in k[x]$ of this
ideal. If we assume $P$ monic, without loss of generality, then $P$ is
uniquely determined.

\begin{definition}
\label{definition-minimal-polynomial}
The polynomial $P$ above is called the {\it minimal polynomial}
of $\alpha$ over $k$.
\end{definition}

\noindent
The minimal polynomial has the following characterization: it is the monic
polynomial, of smallest degree, that annihilates $\alpha$. Any nonconstant
multiple of $P$ will have larger degree, and only multiples of $P$ can
annihilate $\alpha$. This explains the name {\it minimal}.

\medskip\noindent
Clearly the minimal polynomial is {\it irreducible}. This is equivalent to the
assertion that the ideal in $k[x]$ consisting of polynomials annihilating
$\alpha$ is prime. This follows from the fact that the map
$k[x] \to E, x \mapsto \alpha$ is a map into a domain (even a field), so the
kernel is a prime ideal.

\begin{lemma}
\label{lemma-degree-minimal-polynomial}
The degree of the minimal polynomial is $[k(\alpha) : k]$.
\end{lemma}

\begin{proof}
This is just a restatement of the argument in
Lemma \ref{lemma-field-extension-generated-by-one-element}: the observation
is that if $P$ is the minimal polynomial of $\alpha$, then the map
$$
k[x]/(P) \to k(\alpha), \quad x \mapsto \alpha
$$
is an isomorphism as in the aforementioned proof, and we have counted the
degree of such an extension (see
Example \ref{example-degree-simple-algebraic-extension}).
\end{proof}

\noindent
So the observation of the above proof is that if $\alpha \in E$ is algebraic,
then $k(\alpha) \subset E$ is isomorphic to $k[x]/(P)$.


